% Source: sources/returns/S6_04_MAPS_MANIFOLDS/sections/manifolds.tex:5855–6092
\paragraph{Theorem D.2.29 (Proved: all free finite twists in the primitive-$\gamma$ sector).}
The restriction $\gcd(m_i,n_i)=1$ in the homology part of D.2.26
can be removed. The resulting exact formula retains torsion in
the finite cap itself. We keep the same natural initial rows,
minimal original and Mumford fillings, and the finite sites
whose integral fixed lattice contains $\ell_i$ with
$\gamma^*(\ell_i)=1$. These are the $68$ marked finite positions
certified in D.2.26; their orders and indices are $(d_i,q_i)$
from (D.HY2).

At such a position write
\[
 m_i=d_i k_i,\qquad v_i=n_i\ell_i+r_i\delta,\qquad
 t_i=\gcd(m_i,n_i,q_i r_i),\qquad k_i>0.
\]
Every integer power vector and affine lift in this sector is
allowed exactly when
\[
 \gcd(m_i,n_i,r_i)=1,\qquad d_i\mid m_i/t_i.                 \tag{D.UF1}
\]
This is an exact freeness criterion for all powers of the
cyclic action. Under it $t_i$ is either $1$ or $q_i$; the latter
case, when $q_i>1$, is permitted precisely when $q_i\mid k_i$.
The finite central quotient surface $F_i$ and its complex
normal line $L_i$ satisfy
\[
\begin{gathered}
 H_1(F_i)=\mathbb Z^2\oplus\mathbb Z/t_i,\qquad
 H_2(F_i)=\mathbb Z^2\oplus\mathbb Z/t_i,\\
 H_3(F_i)=\mathbb Z^2,\qquad H_0(F_i)=H_4(F_i)=\mathbb Z,\\
 \operatorname{ord}c_1(L_i)=t_i.
\end{gathered}                                               \tag{D.UF2}
\]
Here order one means the zero class. Thus the normal line need
not be topologically trivial.

The full global degree-two presentation has a fibre generator
$s$, two fixed sweeps $D_i,L_i'$ for every finite log, and the
four sweeps $t_j(\Lambda)$ for every regular log. The finite
relations are
\[
\begin{aligned}
 (bn_i/q_i)s+(m_i/q_i)D_i&=0,\\
 -br_i s+m_iL_i'&=0,\\
 r_iD_i+n_iL_i'&=0.
\end{aligned}                                                \tag{D.UF3}
\]
The regular-log relations remain those of (D.SC2), and the
single global relation is
\[
 \sum_iD_i+\sum_jt_j(\delta)+be\,s=0.                       \tag{D.UF4}
\]
There is no coprimality assumption on $m_i,n_i$ beyond the
actual freeness criterion (D.UF1). The regular logs still have
arbitrary free affine vectors.

Let $T$ be the torsion of the cokernel of this matrix. Let
$\mathbb Z^r\oplus A_1$ be the Smith cokernel of the full
Seifert matrix of D.2.23, including these finite numerators
and all regular numerators. Then $r=0$ or $1$, and the complete
integral homology is
\[
\begin{array}{c|c}
 j&H_j(X;\mathbb Z)\\\hline
0,6&\mathbb Z\\
1&\mathbb Z^r\oplus A_1\\
2,3&\mathbb Z^{2r}\oplus T\\
4&\mathbb Z^{2r}\oplus A_1\\
5&\mathbb Z^r.
\end{array}                                                  \tag{D.UF5}
\]
The complete fundamental group is still the Seifert
presentation, rather than its abelianization. In particular
this extends the non-star minimal-filling homology algorithm
to every free affine twist at each of its primitive-$\gamma$
finite sites.

\begin{proof}
\emph{1. Exact freeness and the cap groups.}
Let $A$ be the finite monodromy, of order $d$, and omit the
site index. Its nontrivial eigenvalues are the two eigenvalues
of the corresponding finite elliptic monodromy. Consequently
$A^j$ has the same rational fixed subspace as $A$ whenever
$d\nmid j$. Write $P$ for the rational projection to this fixed
subspace. The integral certificates of D.2.26 give
\[
 \operatorname{Fix}(A)=\mathbb Z\delta\oplus\mathbb Z\ell,
 \qquad P\Lambda=\mathbb Z(\delta/q)\oplus\mathbb Z\ell,
 \qquad N_m\beta=(m/q)\delta,                               \tag{D.UF6}
\]
where $N_m=1+A+\cdots+A^{m-1}$ and $\beta$ is the actual
primitive vertical coinvariant used there. The same certificates
give
$\Lambda/(A-1)\Lambda=\mathbb Z\beta\oplus\mathbb Z\ell$,
with $\delta=q\beta$ in that quotient.

Every affine action with the prescribed power is conjugate
on the good-reduction root disc to the invariant translation
form $h(x)=Ax+v/m$. Indeed its moving translation part has
zero finite-group average and is a translation coboundary;
the averaging solution is holomorphic on the disc. This is
the actual affine comparison of D.2.26, and it does not use
$\gcd(m,n)=1$. Constant unit choices extend over the cap and
preserve the marked normal comparison.

If $g=\gcd(m,n,r)>1$, then $h^{m/g}$ has integral translation
part and fixes the origin of its linear action. Freeness
therefore requires $g=1$. Assume this now. A power with
$d\mid j$ has identity linear part and nonintegral translation,
so is free. For $d\nmid j$, its moving linear map is surjective
on the moving torus. A fixed point exists precisely when its
fixed quotient translation is zero. By (D.UF6) this is
\[
 m\mid jn,\qquad m\mid jqr.
\]
The relevant powers are exactly the nonzero multiples of
$m/t$ below $m$, where $t=\gcd(m,n,qr)$. Since $g=1$ and
$q\in\{1,2,3\}$, $t$ divides $q$. Those powers all have
identity linear part exactly when $d\mid m/t$. This proves
(D.UF1) and the stated alternative for $t$. Freeness on the
central compact torus implies freeness on a sufficiently small
root neighborhood for this finite group, so the criterion is
also sufficient for the smooth holomorphic filling.

The quotient $F=T^4/\langle h\rangle$ has the affine fundamental
group presentation with the extra generator $c$ and the power
relation $mc=v$. On abelianizing, (D.UF6) gives
\[
 H_1(F)=\langle\beta,\ell,c\mid
                   mc-n\ell-qr\beta=0\rangle
        =\mathbb Z^2\oplus\mathbb Z/t.                     \tag{D.UF7}
\]
It is an oriented four-manifold of Euler number zero and
rational first Betti number two. Integral Poincare duality and
the universal coefficient theorem give the other groups in
(D.UF2), including the order-$t$ summand of $H_2$.

The normal line is the associated line for the root rotation
character. On $H_1(F)$ that character sends $c$ to $1/m$ in
$\mathbb R/\mathbb Z$ and the original fibre generators to
zero. The torsion generator
\[
 (m/t)c-(n/t)\ell-(qr/t)\beta
\]
has character $1/t$. The Bockstein of this character is
$c_1(L)$, so its order is exactly $t$. This proves the last
assertion of (D.UF2) without treating a torsion Chern class as
zero.

\emph{2. The full integral Thom kernel.}
The boundary is the marked mapping torus of $A$ and has
$H_1=\mathbb Z^3$, $H_2=\mathbb Z^4$. These statements follow
from the integral local Smith matrices, whose nonzero factors
are all units, in every one of the $68$ positions. The cap is
the disk bundle of $L\to F$. Since $c_1(L)$ is torsion, the
homology Gysin map $H_2(F)\to H_0(F)$ is zero. The pair sequence
therefore makes $H_2(\partial U)\to H_2(U)=H_2(F)$ onto,
including its torsion summand.

The full kernel is the image of the Thom map
\[
 H_3(U,\partial U)=H_1(F)\longrightarrow H_2(\partial U).
\]
The normal circle has the central group representative
$\mu=c^m v^{-1}$. For a lift of each of $\beta,\ell,c$, sweep
that lift by $\mu$. These three normal tori generate the
entire Thom image. On the boundary mapping cylinder, their
normal projections and their fibre components are respectively
\[
\begin{array}{c|c|c}
\text{generator}&\text{normal projection}&
                          \text{fibre component}\\\hline
\beta&N_m\beta&-\beta\wedge v\\
\ell&m\ell&-\ell\wedge v\\
c&v&0.
\end{array}                                                  \tag{D.UF8}
\]
For the first row, subdividing the $m$-fold normal circle into
its $m$ mapping-cylinder intervals gives the sum $N_m\beta$;
the final power-vector edge supplies the displayed fibre
parallelogram. The second row is the same construction for an
invariant loop. The third is the torus of the two commuting
words $c$ and $c^m v^{-1}$, so its only normal sweep is $v$.
These are integral marked cycles, not invariant rational
forms.

Evaluation by the primitive exterior functional $\omega$,
using $-\omega(\beta,\ell)=b/q$ and
$\omega(\ell,\delta)=b$, turns (D.UF8) into the three columns
\[
 \begin{pmatrix}bn/q\\m/q\\0\end{pmatrix},\qquad
 \begin{pmatrix}-br\\0\\m\end{pmatrix},\qquad
 \begin{pmatrix}0\\r\\n\end{pmatrix}                       \tag{D.UF9}
\]
on $(s,D,L')$. The target $H_2(\partial U)$ is free, so the
order-$t$ class of $H_1(F)$ maps to zero. Accordingly the sole
primitive relation among these columns is
\[
 (m/t)K_c-(n/t)K_\ell-(qr/t)K_\beta=0.
\]
No further saturation is allowed or needed: they are images
of generators of the whole Thom source. The projected
lattice has the exact index
\[
 \gcd\left(m^2/q,mn/q,mr\right)=(m/q)t.                    \tag{D.UF10}
\]
In particular a torsion cap multiplies this index by $t$.
When $\gcd(m,n)=1$, the same columns reduce by Bezout to
(D.HY3), recovering the previous theorem.

\emph{3. Global gluing and duality.}
Every degree-two boundary map is onto for these finite caps,
for the regular logs and for the retained minimal original and
Mumford caps. Eliminate their degree-two groups in
Mayer--Vietoris. The exterior fixed-sweep lattice is the
primitive direct sum with its single diagonal $\delta$
relation, as proved for all $172$ rows in D.2.26.
The full kernels (D.UF9), the regular-log kernels and that
one diagonal relation give exactly (D.UF3)--(D.UF4).
This calculation retains the cap torsion through its full
Thom image; it does not assume that the cap or its normal
line has torsion-free invariants.

The rational rank of this cokernel is $r$: after solving the
local sweep relations, the only remaining scalar equation is
$b(e-\sum n_i/m_i)s=0$, with the regular numerators included.
It has the same zero criterion as the Seifert matrix. The
kernel of the degree-one Mayer--Vietoris map is a subgroup
of the free boundary group and is therefore free. Its rank
is also $r$, because the source and target have equal free
rank and their cokernel has free rank $r$. Thus its extension
by the degree-two cokernel splits as an abelian group, giving
$H_2=\mathbb Z^{2r}\oplus T$. All caps still have the same
Euler contribution, so $\chi(X)=2$. Duality and the universal
coefficient theorem now give every remaining entry of (D.UF5).
This includes $e-\sum n_i/m_i=0$ and does not suppress its
free groups or its torsion.
\end{proof}
